
\chapter{REMODEL and REMODEL-V Supplementary Materials}
\label{app:remodel_supplementary}

This appendix provides supplementary
implementation details and illustrative
execution records for the LLM-assisted
architectural design space exploration
frameworks developed in Chapters~2 and~3.

The first section presents selected
LLM prompt specifications that document
structured architectural reasoning,
candidate generation, decision-making,
and adaptive exploration control.
The second section provides a detailed
state-machine execution case study
illustrating candidate processing,
architectural evaluation, and
feedback-guided exploration.

\section{LLM Prompt Specifications}
\label{app:llm_prompt_specifications}

The following prompt excerpts document
selected reasoning and decision-making
instructions used in the REMODEL-V
framework developed in Chapter~3.
These instructions extend the
structured exploration approach
introduced in REMODEL by incorporating
soft-error vulnerability and
feasibility-dependent reasoning
priorities.


\subsection{Cost-Recovery Mode}
\label{app:remodelv_prompts_mode1}


The following excerpt presents the Cost-Recovery Mode
instructions from the final REMODEL-V ACG prompt.

\begin{lstlisting}[
caption={REMODEL-V Cost-Recovery Mode prompt excerpt},
label={lst:remodelv_cost_recovery},
breaklines=true,
basicstyle=\ttfamily\footnotesize,
numbers=left,
numberstyle=\tiny,
frame=single
]
MODE 1 - COST-RECOVERY MODE:

- If REF ranked_executed_designs is empty OR no previous
  design is cost_feasible:
  This means no cost-feasible design has been found yet.

- In this case:
  1) Do not prioritize performance or vulnerability
     improvement yet.
  2) Your ONLY goal is to generate a cost-feasible design.

- Prefer conservative lower-cost settings:
  - Prefer lower-resource or lower-complexity values
    from the provided design_space.
  - Identify which parameters are likely contributing
    to violated cost metrics.
  - Move those parameters toward safer/lower-cost
    values using only allowed design_space values.

IMPORTANT:
- Do NOT repeat a previously executed design.
- Do NOT make only negligible parameter changes.
- If consecutive designs violate the same cost
  constraint, increase the magnitude of the design change.
- Move progressively toward lower-cost regions until
  feasibility is achieved.
- Change at least one parameter that is likely to affect
  the violated cost metric.
- Use previous execution results to identify which
  parameters may be contributing to the violation.
- Avoid getting trapped in repeatedly generating nearly
  identical conservative designs.

- Cost feasibility is the priority, but exploration
  is still required.
\end{lstlisting}


\subsection{Observable-Improvement Mode}
\label{app:remodelv_prompts_mode2}


The following excerpt presents the
Observable-Improvement Mode instructions from
the final REMODEL-V ACG prompt.

\begin{lstlisting}[
caption={REMODEL-V Observable-Improvement Mode prompt instructions},
label={lst:remodelv_observable_improvement},
breaklines=true,
basicstyle=\ttfamily\footnotesize,
numbers=left,
numberstyle=\tiny,
frame=single
]
MODE 2 - OBSERVABLE-IMPROVEMENT MODE:

- If REF contains at least one cost_feasible design:
  This means a cost-feasible design already exists.

- In this case:
  1) Your goal is to improve beyond the top-ranked
     cost-feasible design.
  2) Do NOT only minimize cost.
  3) Do NOT keep repeating the most conservative
     low-cost design.
  4) Use the highest-ranked designs in
     ranked_executed_designs as anchors.
  5) Propose one controlled improvement that may
     improve observable_score while staying
     cost-feasible.
  6) If several consecutive cost-feasible designs
     show similar observable_score, increase
     exploration and consider larger parameter changes.

- Controlled improvements may include:
  - modestly increasing one performance-related
    parameter at a time
  - adjusting one design parameter at a time using
    the provided design_space
  - preserving parameter values that already
    support cost feasibility
  - avoiding aggressive changes that may violate
    cost constraints
\end{lstlisting}


\subsection{Analysis and Candidate Generation}
\label{app:remodelv_prompts_acg}

The following excerpt presents the structured
reasoning and candidate-generation requirements
used in the final REMODEL-V ACG prompt.

\begin{lstlisting}[
caption={REMODEL-V ACG structured reasoning prompt excerpt},
label={lst:remodelv_acg_reasoning},
breaklines=true,
basicstyle=\ttfamily\footnotesize,
numbers=left,
numberstyle=\tiny,
frame=single
]
OUTPUT:
- retrospect: summarize REF context only

- analysis: concise reasoning using REF, ranking,
  cost constraints, performance, and vulnerability
  trade-offs

- analysis.steps must NOT be empty.
- Each analysis step must include one explicit
  parameter-to-metric reasoning claim.
- Each step should identify:
  1) one design parameter,
  2) one metric,
  3) the expected direction
     (increase, decrease, or maintain).

- Example:
  "Increasing <design_parameter> may increase
   <performance_metric>."

- Example:
  "Reducing <design_parameter> may decrease
   <cost_metric>."

- Example:
  "Changing <design_parameter> may affect
   <selected_metric>."

- Avoid vague statements such as:
  * "improve performance"
  * "better design"
  * "stronger configuration"
  unless a specific metric is named.

- suggested_candidates_json: list with exactly
  ONE fully specified design dict

- explanation: brief reason why this candidate
  may improve observable_score while staying
  cost-feasible
\end{lstlisting}


\subsection{Pre-Decision Evaluation}
\label{app:remodelv_prompts_pre}


The following excerpt presents the candidate-assessment
instructions used in the final REMODEL-V PRE prompt.

\begin{lstlisting}[
caption={REMODEL-V PRE candidate-assessment prompt excerpt},
label={lst:remodelv_pre},
breaklines=true,
basicstyle=\ttfamily\footnotesize,
numbers=left,
numberstyle=\tiny,
frame=single
]
YOU ARE IN PRE STATE.
You MUST output a PRE response.

Your task is to evaluate each candidate BEFORE DEC.
PRE does not choose the final candidate.
PRE only decides whether each candidate should
be forwarded to DEC.

INPUTS:
- candidate_pool
- ref_context
- design_space
- user requirements

IMPORTANT RULES:
- Use REF context only.
- Base your reasoning only on observed metrics,
  REF ranking, cost constraints, vulnerability
  metrics, and user requirements.
- Do NOT invent new candidates.
- Do NOT modify candidate parameters.
- Evaluate every candidate in candidate_pool.
- If ranked_executed_designs is empty,
  forward all candidates to DEC.
- If no previous design is cost_feasible,
  forward candidates that are likely to
  improve cost feasibility.
- If at least one previous design is cost_feasible,
  forward candidates that are likely to maintain
  cost feasibility and improve observable_score.

MAIN QUESTION:
For each candidate, ask:
"Should this candidate be forwarded to DEC based
on expected cost feasibility and observable improvement?"

OUTPUT (MANDATORY):
- candidate_assessments: list of one assessment
  per candidate
- Each assessment must include:
  - candidate_id
  - send_to_DEC
  - confidence
  - reason
\end{lstlisting}


\subsection{Decision}
\label{app:remodelv_prompts_dec}


The following excerpt presents the candidate-selection
instructions used in the final REMODEL-V DEC prompt.

\begin{lstlisting}[
caption={REMODEL-V DEC candidate-selection prompt excerpt},
label={lst:remodelv_dec},
breaklines=true,
basicstyle=\ttfamily\footnotesize,
numbers=left,
numberstyle=\tiny,
frame=single
]
YOU ARE IN DEC STATE.
You MUST output a DEC response.

You are the decision engine. Your task is to rank
multiple valid ACG candidates (C0, C1, ..., Cn)
and choose exactly ONE best candidate.

ASSUMPTIONS:
- All candidates are pre-filtered.
- You may ONLY choose from the given candidate IDs.

SELECTION ALGORITHM (MANDATORY):

DEC MUST apply lexicographic reasoning:
1) First, prefer candidates that are most likely
   to remain cost-feasible.
2) Among cost-feasible candidates, prefer stronger
   combined performance and vulnerability behavior.
3) Then prefer lower overall risk.

Step 1 - Per-candidate parameter inspection
For EACH candidate Ci, evaluate every design
parameter listed in the provided design_space.

For each parameter, assess its impact as:
HELP / NEUTRAL / HURT on:
(A) performance objectives
(B) cost objectives
(C) vulnerability objectives

Step 2 - Candidate-level summary
For each Ci, derive:
- predicted_performance_effect
- predicted_vulnerability_effect
- predicted_cost_effect
- combined_observable_effect
- risk_level

Step 3 - Ranking rule (STRICT)
Rank candidates by:
1) better expected cost feasibility
2) stronger combined improvement in performance
   and vulnerability
3) lower overall risk

Performance and vulnerability should be considered
equally important once cost feasibility is satisfied.

Step 4 - Tie-breaking
If tied, choose the candidate whose parameters
most closely follow the strongest history-grounded
conclusions. Prefer coherence over aggressiveness.

Step 5 - Final decision
- Select exactly ONE existing candidate ID.
- Do NOT invent, modify, or merge candidates.

OUTPUT (MANDATORY):
- chosen_id = selected candidate ID
- explanation = brief justification referencing
  the ranking factors above
\end{lstlisting}

\subsection{Insight Extraction}
\label{app:remodelv_prompts_ins}


The following excerpt presents the feasibility-aware
feedback instructions used in the final REMODEL-V
INS prompt.

\begin{lstlisting}[
caption={REMODEL-V INS feasibility-aware feedback prompt excerpt},
label={lst:remodelv_ins_feedback},
breaklines=true,
basicstyle=\ttfamily\footnotesize,
numbers=left,
numberstyle=\tiny,
frame=single
]
MODE 1 - COST FAILURE MODE:

- If the executed design violates one or more
  user-defined cost constraints:
  This means the design is not cost-feasible.

- In this case:
  - Do NOT recommend performance improvements yet.
  - Focus first on restoring cost feasibility.
  - Identify which cost metrics violated or
    approached the bounds when available.
  - Recommend moving toward lower-cost configurations.
  - guidance_for_next_acg must prioritize cost
    feasibility first.

MODE 2 - COST-FEASIBLE OBSERVATION MODE:

- If the executed design satisfies the user-defined
  cost constraints:
  This means the design is cost-feasible.

- In this case:
  - Do NOT over-focus on cost reduction unless a
    cost metric is near its bound.
  - Analyze whether the design provides strong
    performance and vulnerability behavior.
  - Recommend controlled exploration that may
    improve the balance between performance and
    vulnerability while preserving cost feasibility.
  - Avoid repeatedly recommending only conservative
    low-cost settings.
  - Preserve patterns that support cost feasibility
    and strong observed metrics.
\end{lstlisting}




\section{State-Machine Execution Case Study}
\label{app:state_machine_case_study}

This section presents a detailed
execution case study of REMODEL-V
using recorded data from a Matrix
Multiplication (MM) experiment
under the Medium cost-constraint
regime.

The case study illustrates how
user-defined optimization requirements
are processed through the REMODEL-V
state machine, from architectural
candidate generation and deterministic
verification to execution and
feedback-guided exploration.

The selected experiment uses six
parallel Architectural Candidate
Generation (ACG) instances and
a maximum exploration budget
of five iterations.
Two consecutive iterations
are examined to illustrate
the transition from
Cost-Recovery Mode to
Observable-Improvement Mode.

The walkthrough combines recorded
LLM responses and architectural
measurements with explicitly
constructed candidate-processing
examples. The constructed examples
illustrate framework behavior
and are not presented as
additional experimental observations.


\subsection{User Requirements and Optimization Constraints}
\label{app:case_study_requirements}

The execution begins with the user specifying
the architectural design space, optimization
objectives, hard cost constraints, and exploration
requirements.

The selected experiment uses the MM workload
with database-backed architectural evaluation.
In this execution mode, the framework retrieves
previously measured architectural outcomes for
selected configurations rather than launching
new simulations during exploration.

The user selects four performance objectives
to maximize: IPC, cache hit rate, MSHR hit rate,
and access throughput. Two soft-error vulnerability
objectives are minimized: L2-cache MSHR vulnerability
per tick and L1D-cache MSHR vulnerability per tick.

The optimization follows a cost-sensitive mission
in which all seven cost constraints must be
satisfied before a configuration is considered
cost-feasible.

\subsubsection{Architectural Parameter Domains}

The user-defined architectural design space
contains six configurable parameters.
Table~\ref{tab:appendix_design_space} summarizes
their permitted values.

\begin{table}[htbp]
\centering
\caption{Architectural parameter domains supplied
to REMODEL-V in the selected execution.}
\label{tab:appendix_design_space}
\renewcommand{\arraystretch}{1.1}
\begin{tabular}{ll}
\hline
\textbf{Parameter} & \textbf{Allowed Values} \\
\hline
L1D size & 8, 16, 32, 64 kB \\
L1D MSHRs & 4, 6, 12 \\
ROB entries & 64, 128, 192, 256 \\
LQ entries & 8, 16, 32 \\
L2 size & 128, 256, 512 kB; 1, 2 MB \\
L2 associativity & 2, 4, 8, 16, 32 \\
\hline
\end{tabular}
\end{table}

These domains define 3,600 possible architectural
configurations. Every generated candidate must
assign a valid value to each parameter.

The parameter domains are subsequently used by
the deterministic Feasibility Verification (FV)
state to reject configurations containing
missing, unknown, or out-of-domain parameter
values.

\subsubsection{Hard Cost Constraints}

The user specifies seven hard cost constraints
under the Medium regime, as summarized in
Table~\ref{tab:appendix_cost_bounds}.

\begin{table}[htbp]
\centering
\caption{Hard cost constraints used in the
selected Medium-regime execution.}
\label{tab:appendix_cost_bounds}
\renewcommand{\arraystretch}{1.1}
\begin{tabular}{lcc}
\hline
\textbf{Cost Metric} &
\textbf{Lower Bound} &
\textbf{Upper Bound} \\
\hline
Total power & 0 & 2.5 \\
Area & 0 & 5.0 \\
Dynamic power & 0 & 2.3 \\
Peak power & 0 & 9.0 \\
Leakage power & 0 & 0.15 \\
Total energy & 0 & 0.15 \\
EDP & 0 & 8.0 \\
\hline
\end{tabular}
\end{table}

Unlike architectural parameter validity,
cost feasibility depends on measured
architectural outcomes. A configuration is
classified as cost-feasible only when all
seven measured cost values satisfy their
corresponding bounds.

These requirements are provided to the
framework and establish the constraints
under which architectural candidates are
generated, assessed, selected, and evaluated.


\subsection{Iteration 1: Initial Candidate Exploration}
\label{app:illustrative_iteration1}

\subsubsection{REF: Initial Reference Context}
\label{app:illustrative_ref}

The first iteration begins with the Reference
Context (REF) state, which constructs the
information available to the subsequent
reasoning and candidate-generation states.

At initialization, no architectural
configuration has been executed. Consequently,
REF contains no measured architectural
outcomes, previously executed designs, or
execution-grounded insights.

Listing~\ref{lst:illustrative_ref}
presents the initial REF context, adapted
from the recorded REMODEL-V execution.

\begin{lstlisting}[
    style=remodeljson,
    caption={Initial REF context used in the
    illustrative state-machine walkthrough.},
    label={lst:illustrative_ref}
]
{
  "iteration": 1,
  "ranked_executed_designs": [],
  "previously_executed_designs": [],
  "remaining_weaknesses": [],
  "latest_insight": null,
  "guidance_for_next_acg": [
    "Use only observed metrics and constraints from REF context.",
    "Prioritize cost feasibility first.",
    "Among cost-feasible designs, prefer higher observable_score.",
    "Do not repeat previously executed designs."
  ]
}
\end{lstlisting}

The REF output is interpreted as follows:

\begin{enumerate}

\item \textbf{Iteration identification (Line 2):}
Identifies the beginning of the first
exploration iteration.

\item \textbf{Execution history (Lines 3--4):}
No previously executed configurations
are available for ranking or comparison.

\item \textbf{Execution-grounded feedback
(Lines 5--6):}
No previous weaknesses or insights
are available because no configuration
has been evaluated.

\item \textbf{Candidate-generation guidance
(Lines 7--12):}
REF instructs ACG to prioritize cost
feasibility, consider the selected
performance and vulnerability objectives,
and avoid repeating previously executed
configurations.

\end{enumerate}

Because no cost-feasible configuration
has yet been observed, Cost-Recovery Mode
is active. The resulting REF context,
together with the user-defined requirements,
is supplied to the six independent ACG
instances.


\subsubsection{ACG: Architectural Candidate Generation}
\label{app:illustrative_acg}

Following REF context construction,
REMODEL-V enters the Architectural
Candidate Generation (ACG) state.

In this illustrative scenario, six
independent ACG instances operate in
parallel. Each instance receives the
same user requirements, architectural
parameter domains, and REF context,
but independently generates one
architectural proposal.

The candidate-generation process is
guided by Cost-Recovery Mode, which
prioritizes establishing cost feasibility
while considering the selected performance
and vulnerability objectives.

Table~\ref{tab:illustrative_acg_candidates}
presents the six candidate configurations
used in this worked example.

\begin{table*}[htbp]
\centering
\caption{Six architectural candidates used
in the illustrative ACG example.}
\label{tab:illustrative_acg_candidates}
\small
\renewcommand{\arraystretch}{1.1}
\begin{tabular}{lcccccc}
\hline
\textbf{Parameter} &
\textbf{ACG\_0} &
\textbf{ACG\_1} &
\textbf{ACG\_2} &
\textbf{ACG\_3} &
\textbf{ACG\_4} &
\textbf{ACG\_5} \\
\hline
L1D size & 8 & 16 & 32 & 16 & 16 & 64 \\
L1D MSHRs & 4 & 6 & 6 & 5 & 6 & 12 \\
ROB entries & 64 & 128 & 128 & 128 & 128 & 256 \\
LQ entries & 8 & 16 & 16 & 16 & 16 & 32 \\
L2 size & 128 & 256 & 256 & 256 & 256 & 1024 \\
L2 associativity & 2 & 4 & 4 & 4 & 4 & 8 \\
\hline
\end{tabular}

\vspace{2pt}
\footnotesize
L1D and L2 capacities are expressed in kB.
\end{table*}

The first three configurations reproduce
candidate proposals from a recorded
REMODEL-V execution. The remaining
three candidates are constructed to
illustrate additional candidate-processing
behaviors.

Specifically, ACG\_3 contains an
out-of-domain L1D MSHR value, ACG\_4
duplicates ACG\_1, and ACG\_5 provides
an additional valid configuration.

These illustrative proposals are used
to demonstrate the behavior of FV and
CDR/DED. They are not presented as
additional experimental observations.



\paragraph{Recorded LLM Response for ACG\_1}
\label{app:illustrative_acg_reasoning}

To illustrate the reasoning process within
Architectural Candidate Generation (ACG),
Listing~\ref{lst:illustrative_acg1}
presents the recorded LLM response for ACG\_1.

The example preserves the architectural
parameter assignment and structured
reasoning generated by the LLM during
an actual REMODEL-V execution.

\begin{lstlisting}[
    style=remodeljson,
    caption={Recorded ACG\_1 architectural proposal
    and structured LLM reasoning.},
    label={lst:illustrative_acg1}
]
{
  "candidate_id": "ACG_1",
  "params": {
    "l1d_size": "16kB",
    "l1d_mshrs": 6,
    "rob_entries": 128,
    "lq_entries": 16,
    "l2_size": "256kB",
    "l2_assoc": 4
  },
  "trace": {
    "retrospect": "No previous designs have been executed.",
    "analysis": {
      "steps": [
        {
          "thinking": "Choosing l1d_size of '16kB' provides a balance between area and access speed, maintaining acceptable performance metrics.",
          "step_conclusion": "Increasing l1d_size can enhance hit_rate."
        },
        {
          "thinking": "Setting l1d_mshrs to 6 allows for a moderate increase in cache efficiency without significantly increasing cost.",
          "step_conclusion": "Increasing l1d_mshrs may improve ipc."
        },
        {
          "thinking": "Using 128 rob_entries optimizes the balance between complexity and performance, ensuring an increase in task scheduling efficiency.",
          "step_conclusion": "Increasing rob_entries may lead to better access_throughput."
        },
        {
          "thinking": "Setting lq_entries to 16 provides a robust queue length supporting the enhanced throughput and latency management.",
          "step_conclusion": "Increasing lq_entries would improve access_throughput."
        },
        {
          "thinking": "Choosing l2_size of '256kB' ensures sufficient capacity for a cache while remaining within the cost boundaries.",
          "step_conclusion": "Increasing l2_size may lead to better hit_rate."
        },
        {
          "thinking": "Setting l2_assoc to 4 strikes a balance between complexity and hit rate, aiding in vulnerability reduction.",
          "step_conclusion": "Increasing l2_assoc may decrease system vulnerability."
        }
      ],
      "final_conclusion": "This design maximizes performance metrics while keeping cost within user-defined constraints."
    },
    "explanation": "This candidate is structured to ensure cost feasibility while aiming to enhance observable performance metrics."
  }
}
\end{lstlisting}

\paragraph{Interpretation of the ACG Response}

Listing~\ref{lst:illustrative_acg1}
demonstrates how ACG produces both a
candidate configuration and a structured
reasoning trace.

The response contains five principal
elements:

\begin{enumerate}

\item \textbf{Candidate identification and
architectural parameters (Lines 2--10):}
The response identifies ACG\_1 and assigns
values to all six architectural parameters.
These values form the proposed configuration
that is subsequently processed by NOR and FV.

\item \textbf{Retrospective context (Line 13):}
The \texttt{retrospect} field states that
no previous configurations have been executed.
The LLM therefore begins without measured
architectural evidence from earlier iterations.

\item \textbf{Parameter-level reasoning
(Lines 14--40):}
The \texttt{analysis.steps} array contains
six reasoning entries corresponding to
the six architectural parameters.

Each entry provides a \texttt{thinking}
statement describing the expected effect
of the selected parameter and a
\texttt{step\_conclusion} summarizing
the architectural hypothesis.

For example, the LLM expects the selected
L1D capacity to support cache hit rate,
the MSHR allocation to influence IPC,
and the ROB capacity to affect
access throughput.

\item \textbf{Overall reasoning conclusion
(Line 41):}
The \texttt{final\_conclusion} summarizes
the LLM's expectation that the configuration
will provide favorable performance while
respecting the active cost constraints.

\item \textbf{Final explanation (Line 43):}
The \texttt{explanation} field provides
a concise justification for the complete
architectural proposal.

\end{enumerate}

These statements are LLM-generated
architectural hypotheses rather than
verified simulation findings. Although
ACG\_1 predicts favorable performance
and cost feasibility, the actual outcomes
remain unknown until the selected design
is evaluated.

The six illustrative candidates are
subsequently considered in the
candidate-processing walkthrough,
which demonstrates normalization,
architectural validity checking,
and duplicate elimination.


\subsubsection{NOR: Candidate Normalization}
\label{app:illustrative_nor}

Following architectural candidate generation,
the framework enters the Normalization (NOR)
state. NOR is a deterministic, code-based
mechanism responsible for converting generated
architectural proposals into a consistent
internal representation.

Each candidate is expected to contain a
complete architectural parameter assignment.
NOR extracts the parameter dictionary and
preserves the candidate identifier and
associated reasoning information.

Unlike ACG, NOR does not generate new
configurations or invoke an LLM. It also
does not determine whether the proposed
parameter values belong to the permitted
architectural domains. That responsibility
belongs to the subsequent FV state.

\paragraph{Illustrative NOR Input and Output}

Listing~\ref{lst:illustrative_nor}
shows a compact representation of the
six illustrative candidates after
normalization. Only candidate identifiers
and architectural parameters are displayed;
auxiliary reasoning fields are omitted.

\begin{lstlisting}[
    style=remodeljson,
    caption={Illustrative normalized candidate
    representations after NOR.},
    label={lst:illustrative_nor}
]
[
  {
    "candidate_id": "ACG_0",
    "params": {
      "l1d_size": "8kB",
      "l1d_mshrs": 4,
      "rob_entries": 64,
      "lq_entries": 8,
      "l2_size": "128kB",
      "l2_assoc": 2
    }
  },
  {
    "candidate_id": "ACG_1",
    "params": {
      "l1d_size": "16kB",
      "l1d_mshrs": 6,
      "rob_entries": 128,
      "lq_entries": 16,
      "l2_size": "256kB",
      "l2_assoc": 4
    }
  },
  {
    "candidate_id": "ACG_2",
    "params": {
      "l1d_size": "32kB",
      "l1d_mshrs": 6,
      "rob_entries": 128,
      "lq_entries": 16,
      "l2_size": "256kB",
      "l2_assoc": 4
    }
  },
  {
    "candidate_id": "ACG_3",
    "params": {
      "l1d_size": "16kB",
      "l1d_mshrs": 5,
      "rob_entries": 128,
      "lq_entries": 16,
      "l2_size": "256kB",
      "l2_assoc": 4
    }
  },
  {
    "candidate_id": "ACG_4",
    "params": {
      "l1d_size": "16kB",
      "l1d_mshrs": 6,
      "rob_entries": 128,
      "lq_entries": 16,
      "l2_size": "256kB",
      "l2_assoc": 4
    }
  },
  {
    "candidate_id": "ACG_5",
    "params": {
      "l1d_size": "64kB",
      "l1d_mshrs": 12,
      "rob_entries": 256,
      "lq_entries": 32,
      "l2_size": "1MB",
      "l2_assoc": 8
    }
  }
]
\end{lstlisting}

\paragraph{Interpretation of NOR Processing}

The normalized candidate representations
illustrate three important properties
of the NOR state.

First, all six candidates are represented
using the same parameter names and dictionary
structure, allowing subsequent states to
process them consistently.

Second, NOR preserves the original
architectural parameter assignments.
For example, ACG\_3 retains an L1D MSHR
count of 5, even though this value is
outside the permitted architectural domain.
NOR does not modify this value because
architectural validity checking is performed
separately by FV.

Third, ACG\_4 remains identical to ACG\_1.
NOR does not eliminate duplicate
configurations because duplicate detection
is handled by CDR/DED.

Consequently, all six illustrative
candidates remain after normalization:

\begin{equation}
N_{\mathrm{ACG}} =
N_{\mathrm{NOR}} = 6.
\end{equation}

The normalized candidate representations
are then forwarded to FV for deterministic
architectural validity checking.


\subsubsection{FV: Feasibility Verification}
\label{app:illustrative_fv}

Following normalization, the candidate pool enters
the Feasibility Verification (FV) state. FV is a
deterministic, code-based mechanism that verifies
whether each architectural configuration satisfies
the structural requirements of the user-defined
design space.

For each candidate, FV checks that all required
architectural parameters are present, no unknown
parameters are introduced, and every assigned value
belongs to its corresponding permitted domain.

This verification does not require an LLM call or
architectural simulation. It also differs from
cost-feasibility assessment, which depends on
measured power, area, energy, and related outcomes.

\paragraph{Illustrative FV Validation Results}

Listing~\ref{lst:illustrative_fv} presents the
validation outcomes for the six candidates
introduced in the preceding ACG and NOR examples.

\begin{lstlisting}[
    style=remodeljson,
    caption={Illustrative FV validation results
    for the six architectural candidates.},
    label={lst:illustrative_fv}
]
[FV] Feasibility verification

ACG_0: ACCEPTED
  All architectural parameters are valid.

ACG_1: ACCEPTED
  All architectural parameters are valid.

ACG_2: ACCEPTED
  All architectural parameters are valid.

ACG_3: REJECTED
  Invalid value for l1d_mshrs: 5
  Allowed values: [4, 6, 12]

ACG_4: ACCEPTED
  All architectural parameters are valid.

ACG_5: ACCEPTED
  All architectural parameters are valid.

[FV] Accepted candidates: 5
[FV] Rejected candidates: 1
\end{lstlisting}

\paragraph{Interpretation of FV Results}

The illustrative validation results demonstrate
how FV separates structurally valid architectural
configurations from invalid proposals.

\begin{enumerate}

\item \textbf{Accepted candidates (Lines 3--11
and 16--20):}
ACG\_0, ACG\_1, ACG\_2, ACG\_4, and ACG\_5
contain permitted values for all six
architectural parameters. They therefore
pass deterministic validation.

\item \textbf{Rejected candidate (Lines 12--14):}
ACG\_3 specifies an L1D MSHR count of 5.
However, the user-defined parameter domain is

\begin{equation}
\mathcal{M}_{\mathrm{L1D}} = \{4,6,12\}.
\end{equation}

Since

\begin{equation}
5 \notin \mathcal{M}_{\mathrm{L1D}},
\end{equation}

the candidate violates the architectural
parameter constraints and is rejected.

\item \textbf{Remaining candidate pool
(Lines 22--23):}
Five candidates remain eligible for
subsequent processing. The invalid
configuration is excluded before
candidate selection or architectural
evaluation.

\end{enumerate}

The candidate count changes from six
normalized proposals to five valid
configurations:

\begin{equation}
N_{\mathrm{NOR}} = 6,
\qquad
N_{\mathrm{FV}} = 5.
\end{equation}

Importantly, FV does not determine whether
the accepted configurations satisfy the
user-defined hardware cost bounds.
It verifies architectural legality only.

Furthermore, ACG\_4 is accepted even though
its configuration duplicates ACG\_1.
Duplicate detection is not the responsibility
of FV and is instead performed by the
subsequent CDR/DED state.

The five accepted configurations are
therefore forwarded to CDR/DED for
canonicalization and duplicate elimination.


\subsubsection{CDR/DED: Canonicalization and Deduplication}
\label{app:illustrative_cdr}

Following architectural validity verification,
the five accepted candidates enter the
Canonicalization and Deduplication (CDR/DED)
state.

CDR/DED is a deterministic, code-based
mechanism that prevents redundant architectural
configurations from proceeding to subsequent
decision-making and execution states.

For each candidate, the framework constructs
a canonical representation of its architectural
parameter assignment. This representation
is used to identify configurations that
are duplicated within the current candidate
pool or have already been executed in
previous iterations.

In the implementation, the canonical key
is constructed by sorting the architectural
parameter--value pairs:

\begin{equation}
K(x) =
\operatorname{tuple}
\left(
\operatorname{sort}
\left(
\{(p_i,v_i)\}_{i=1}^{6}
\right)
\right).
\label{eq:illustrative_canonical_key}
\end{equation}

Two candidates are considered duplicates
when their canonical keys are identical,
regardless of the original ordering of
their parameter assignments.

\paragraph{Illustrative Duplicate Detection}

Listing~\ref{lst:illustrative_cdr}
shows how CDR/DED processes the five
candidates accepted by FV.

The listing is a constructed representation
of the deterministic processing results,
following the duplicate-detection rules
implemented in REMODEL-V.

\begin{lstlisting}[
    style=remodeljson,
    caption={Illustrative CDR/DED processing
    and duplicate elimination.},
    label={lst:illustrative_cdr}
]
[CDR/DED] Canonicalization and deduplication

ACG_0: KEPT
  Unique architectural configuration.
  Assigned ID: C0

ACG_1: KEPT
  Unique architectural configuration.
  Assigned ID: C1

ACG_2: KEPT
  Unique architectural configuration.
  Assigned ID: C2

ACG_4: REJECTED
  Reason: Duplicate in current candidate pool.
  Identical configuration to ACG_1.

ACG_5: KEPT
  Unique architectural configuration.
  Assigned ID: C3

[CDR/DED] Input candidates: 5
[CDR/DED] Duplicates removed: 1
[CDR/DED] Unique candidates retained: 4
\end{lstlisting}

\paragraph{Interpretation of CDR/DED Results}

Listing~\ref{lst:illustrative_cdr}
demonstrates how the framework identifies
and eliminates redundant proposals.

\begin{enumerate}

\item \textbf{Unique candidates (Lines 3--14):}
ACG\_0, ACG\_1, and ACG\_2 have distinct
architectural parameter assignments.
Their canonical keys are therefore
different, and all three are retained.

\item \textbf{Duplicate rejection (Lines 15--17):}
ACG\_4 has exactly the same parameter
assignment as ACG\_1:

\begin{equation}
K(x_{\mathrm{ACG}_4})
=
K(x_{\mathrm{ACG}_1}).
\end{equation}

Because ACG\_1 was already retained
in the current candidate pool,
ACG\_4 is rejected as a duplicate.

\item \textbf{Additional unique candidate
(Lines 19--21):}
ACG\_5 contains a distinct architectural
configuration and is retained.

\item \textbf{Final candidate pool
(Lines 23--25):}
Four unique configurations remain after
duplicate elimination.

\end{enumerate}

The candidate count is reduced from
five valid configurations to four
unique configurations:

\begin{equation}
N_{\mathrm{FV}} = 5,
\qquad
N_{\mathrm{CDR/DED}} = 4.
\end{equation}

The retained candidates are assigned
new internal identifiers, C0 through C3,
while preserving their original ACG
identifiers for traceability.

\begin{table}[htbp]
\centering
\caption{Candidate identifiers after
illustrative CDR/DED processing.}
\label{tab:illustrative_cdr_mapping}
\renewcommand{\arraystretch}{1.1}
\begin{tabular}{lcc}
\hline
\textbf{ACG Candidate} &
\textbf{Internal ID} &
\textbf{Status} \\
\hline
ACG\_0 & C0 & Retained \\
ACG\_1 & C1 & Retained \\
ACG\_2 & C2 & Retained \\
ACG\_3 & -- & Rejected by FV \\
ACG\_4 & -- & Duplicate of ACG\_1 \\
ACG\_5 & C3 & Retained \\
\hline
\end{tabular}
\end{table}

By removing redundant configurations
before candidate assessment and selection,
CDR/DED prevents duplicate proposals
from consuming additional exploration
resources.

The four unique candidates are then
forwarded to PRE for pre-decision assessment.


\subsubsection{PRE: Pre-Decision Evaluation}
\label{app:illustrative_pre}

Following canonicalization and deduplication,
the four remaining architectural candidates
enter the Pre-Decision Evaluation (PRE) state.

PRE assesses candidate suitability before
final selection by DEC. When previous execution
evidence is available, PRE uses LLM-based
reasoning over the REF context, measured
architectural outcomes, and active cost
constraints to determine which candidates
should proceed to DEC.

However, the framework implements a special
initialization rule for the first iteration.
When REF contains no previously executed
designs, PRE does not invoke the LLM for
candidate assessment. Instead, all remaining
candidates are forwarded to DEC because
no measured reference design is available
for comparison.

\paragraph{Illustrative PRE Output}

Listing~\ref{lst:illustrative_pre}
shows the PRE assessments that would be
produced for the four remaining candidates
under this initialization rule.

The listing is constructed using the
deterministic PRE behavior implemented
in REMODEL-V.

\begin{lstlisting}[
    style=remodeljson,
    caption={Illustrative PRE output for the
    four candidates in iteration 1.},
    label={lst:illustrative_pre}
]
[
  {
    "candidate_id": "C0",
    "source_candidate_id": "ACG_0",
    "likely_improves_observed_performance": true,
    "send_to_DEC": true,
    "confidence": 1.0,
    "reason": "No current best design exists yet, so candidate is forwarded to DEC."
  },
  {
    "candidate_id": "C1",
    "source_candidate_id": "ACG_1",
    "likely_improves_observed_performance": true,
    "send_to_DEC": true,
    "confidence": 1.0,
    "reason": "No current best design exists yet, so candidate is forwarded to DEC."
  },
  {
    "candidate_id": "C2",
    "source_candidate_id": "ACG_2",
    "likely_improves_observed_performance": true,
    "send_to_DEC": true,
    "confidence": 1.0,
    "reason": "No current best design exists yet, so candidate is forwarded to DEC."
  },
  {
    "candidate_id": "C3",
    "source_candidate_id": "ACG_5",
    "likely_improves_observed_performance": true,
    "send_to_DEC": true,
    "confidence": 1.0,
    "reason": "No current best design exists yet, so candidate is forwarded to DEC."
  }
]
\end{lstlisting}

\paragraph{Interpretation of PRE Processing}

Listing~\ref{lst:illustrative_pre}
demonstrates the initialization behavior
of PRE.

\begin{enumerate}

\item \textbf{Candidate assessments:}
Each candidate receives an assessment
containing its internal identifier,
original ACG identifier, forwarding
decision, confidence value, and
corresponding explanation.

\item \textbf{Forwarding decisions:}
The \texttt{send\_to\_DEC} field is
set to \texttt{true} for all four
candidates. Therefore, no candidate
is removed during PRE.

\item \textbf{Confidence values:}
Each candidate receives a confidence
value of 1.0. This is a predefined
initialization value rather than an
LLM-estimated probability of architectural
feasibility or performance improvement.

\item \textbf{Assessment explanation:}
The \texttt{reason} field indicates
that no current best design exists.
Consequently, PRE forwards the
candidate without comparing it
against previously measured outcomes.

\end{enumerate}

Thus, the candidate count remains unchanged:

\begin{equation}
N_{\mathrm{CDR/DED}}
=
N_{\mathrm{PRE}}
=
4.
\end{equation}

This behavior ensures that the framework
can perform initial architectural exploration
even when no execution-grounded evidence
is available.

In subsequent iterations, after at least
one design has been executed, PRE can
use the REF context and measured outcomes
to assess candidate quality and cost risk.
It may then exclude candidates that are
unlikely to improve the observed objectives
or preserve cost feasibility.

The four PRE-approved candidates are
forwarded to DEC for final selection.


\subsubsection{DEC: Final Candidate Selection}
\label{app:illustrative_dec}

Following PRE assessment, the framework enters
the Decision (DEC) state, which selects one
architectural configuration from the
PRE-approved candidate pool for execution.

Unlike the preceding deterministic candidate
processing stages, DEC uses LLM-based reasoning
when multiple candidates remain available.

DEC receives the candidate pool, the
corresponding PRE assessments, the REF context,
the user-defined optimization requirements,
and the available execution history.

Its task is to select exactly one candidate
while considering architectural performance,
soft-error vulnerability, cost constraints,
and the active reasoning priority.

\paragraph{Illustrative DEC Input}

In the constructed example, all four
candidates pass PRE and are forwarded to DEC.

Listing~\ref{lst:illustrative_dec_input}
shows a compact representation of the
candidate pool supplied to DEC.

\begin{lstlisting}[
    style=remodeljson,
    caption={Illustrative candidate pool
    supplied to DEC in iteration 1.},
    label={lst:illustrative_dec_input}
]
{
  "current_state": "DEC",
  "candidate_pool": [
    {
      "id": "C0",
      "source": "ACG_0"
    },
    {
      "id": "C1",
      "source": "ACG_1"
    },
    {
      "id": "C2",
      "source": "ACG_2"
    },
    {
      "id": "C3",
      "source": "ACG_5"
    }
  ],
  "history": []
}
\end{lstlisting}

For readability, Listing~\ref{lst:illustrative_dec_input}
shows only the candidate identifiers and their
original ACG sources. In the actual framework,
DEC receives the architectural parameter
assignments, PRE assessments, REF context,
and user-defined requirements.

Because this is the first iteration, no
previously executed configurations are
available in the execution history.

\paragraph{Illustrative DEC Selection}

For this worked example, DEC selects C1,
corresponding to ACG\_1.

Listing~\ref{lst:illustrative_dec_output}
presents the selection using the original
DEC explanation recorded for C1 in a
separate three-candidate REMODEL-V execution.

The explanation is reproduced to illustrate
the type of reasoning generated by DEC;
the four-candidate selection itself is
constructed for this walkthrough.

\begin{lstlisting}[
    style=remodeljson,
    caption={Illustrative DEC selection using
    a recorded DEC explanation for C1.},
    label={lst:illustrative_dec_output}
]
{
  "chosen_id": "C1",
  "explanation": "Candidate C1 features parameters that provide better resource allocation with increased l1d_size and l2_size, likely improving overall performance while remaining cost-feasible."
}
\end{lstlisting}

\paragraph{Interpretation of the DEC Response}

Listing~\ref{lst:illustrative_dec_output}
contains two principal fields.

\begin{enumerate}

\item \textbf{Selected candidate (Line 2):}
The \texttt{chosen\_id} field identifies
C1 as the configuration selected for
architectural evaluation.

DEC selects an existing candidate from
the PRE-approved pool rather than
generating a new configuration.

\item \textbf{Selection explanation (Line 3):}
The \texttt{explanation} field provides
the LLM's justification for selecting C1.

The recorded explanation emphasizes
architectural resource allocation,
particularly L1D and L2 cache capacity,
as potential contributors to favorable
performance while respecting cost
constraints.

\end{enumerate}

The selected candidate has the following
architectural parameters:

\begin{equation}
x_{\mathrm{DEC}}
=
\left(
16\text{ kB},
6,
128,
16,
256\text{ kB},
4
\right),
\label{eq:illustrative_dec_design}
\end{equation}

where the parameter ordering is L1D size,
L1D MSHRs, ROB entries, LQ entries,
L2 size, and L2 associativity.

Importantly, the DEC explanation represents
an architectural expectation rather than
a verified performance or cost result.
At this stage, the selected configuration
has passed architectural validity and
duplicate checks, but its measured
cost feasibility has not yet been established
within the illustrative exploration.

DEC forwards the selected configuration
to EXC for architectural evaluation.



\subsubsection{EXC: Architectural Execution and Measurement}
\label{app:illustrative_exc}

Following candidate selection, the framework enters
the Execution (EXC) state. EXC obtains the
architectural measurements of the configuration
selected by DEC, including performance, hardware
cost, and soft-error vulnerability outcomes.

Unlike the LLM-driven reasoning states, EXC does
not predict architectural behavior. Instead, it
retrieves or executes the selected configuration
using the configured architectural evaluation
backend.

In the original Matrix Multiplication experiment,
the framework operates in database-backed execution
mode. Consequently, the selected configuration is
matched against previously characterized simulation
records, and its measured outcomes are retrieved
rather than launching a new gem5 simulation.

For this illustrative walkthrough, the selected
candidate C1 corresponds to an architectural
configuration evaluated in the recorded
Medium-regime experiment. Its actual measurements
are reused to demonstrate the EXC operation.

\paragraph{Selected Architectural Configuration}

The configuration selected by DEC is

\begin{equation}
x_{\mathrm{C1}} =
(
16\text{ kB},
6,
128,
16,
256\text{ kB},
4
),
\label{eq:illustrative_exc_design}
\end{equation}

where the parameter ordering is L1D size,
L1D MSHRs, ROB entries, LQ entries,
L2 size, and L2 associativity.

\paragraph{Recorded Execution Output}

Listing~\ref{lst:illustrative_exc}
presents the measured performance,
vulnerability, and cost values retrieved
for the selected configuration.

The listing is a condensed representation
of the recorded \texttt{execution\_result}
and preserves the original measured values
for the selected metrics.

\begin{lstlisting}[
    style=remodeljson,
    caption={Recorded architectural measurements
    for the selected candidate C1.},
    label={lst:illustrative_exc}
]
{
  "status": "evaluated_from_db",
  "design": {
    "l1d_size": "16kB",
    "l1d_mshrs": 6,
    "rob_entries": 128,
    "lq_entries": 16,
    "l2_size": "256kB",
    "l2_assoc": 4
  },
  "metrics": {
    "ipc": 1.15112703986,
    "hit_rate": 0.999986879309,
    "mshr_hit_rate": 0.997260273973,
    "access_throughput": 326404710.921,
    "system.l2.mshr.vulnerability_per_tick": 0.004863458035739129,
    "system.cpu.dcache.mshr.vulnerability_per_tick": 0.09580629212425008,
    "total_power": 2.34243826,
    "area": 4.3912,
    "dynamic_power": 2.27206,
    "peak_power": 8.76435826,
    "leakage_power": 0.07037826,
    "total_energy": 0.00109391866742,
    "edp": 5.10860017685e-07
  }
}
\end{lstlisting}

\paragraph{Interpretation of Execution Results}

Listing~\ref{lst:illustrative_exc}
demonstrates how EXC provides measured
architectural evidence for subsequent
reasoning and decision-making.

\begin{enumerate}

\item \textbf{Execution status:}
The \texttt{evaluated\_from\_db} status
indicates that the selected configuration
was matched against a previously
characterized architectural record.

\item \textbf{Architectural configuration:}
The \texttt{design} field identifies the
six architectural parameters of the
selected candidate. These values match
the configuration forwarded by DEC.

\item \textbf{Performance measurements:}
The recorded performance includes an
IPC of approximately 1.1511, a cache
hit rate of 0.99999, an MSHR hit rate
of 0.99726, and an access throughput
of approximately $3.264 \times 10^8$.

\item \textbf{Soft-error vulnerability:}
The measured L2-MSHR vulnerability is
approximately 0.004863 per tick, while
the L1D-MSHR vulnerability is approximately
0.095806 per tick.

These quantities represent measured
vulnerability outcomes rather than
predictions generated by the LLM.

\item \textbf{Hardware cost measurements:}
EXC also retrieves the seven cost metrics
required to determine whether the selected
configuration satisfies the active
Medium-regime constraints.

\end{enumerate}

\paragraph{Cost-Feasibility Verification}

The measured cost values are compared
against the user-defined Medium-regime
bounds introduced in
Table~\ref{tab:appendix_cost_bounds}.

\begin{table}[htbp]
\centering
\caption{Measured cost feasibility of the
selected architectural configuration C1.}
\label{tab:illustrative_exc_feasibility}
\renewcommand{\arraystretch}{1.1}
\small
\begin{tabular}{lccc}
\hline
\textbf{Cost Metric} &
\textbf{Measured} &
\textbf{Upper Bound} &
\textbf{Status} \\
\hline
Total power & 2.3424 & 2.5 & Feasible \\
Area & 4.3912 & 5.0 & Feasible \\
Dynamic power & 2.2721 & 2.3 & Feasible \\
Peak power & 8.7644 & 9.0 & Feasible \\
Leakage power & 0.07038 & 0.15 & Feasible \\
Total energy & 0.001094 & 0.15 & Feasible \\
EDP & $5.109\times10^{-7}$ & 8.0 & Feasible \\
\hline
\end{tabular}
\end{table}

All seven measured cost values satisfy
their corresponding bounds. Therefore,
the selected configuration is classified
as cost-feasible.

Formally,

\begin{equation}
C_j(x_{\mathrm{C1}})
\leq C_{j,\mathrm{Medium}}^{\max},
\qquad \forall j \in \{1,\ldots,7\}.
\end{equation}

Although C1 satisfies all cost constraints,
its dynamic power is close to the maximum
permitted value of 2.3 W. This limited
remaining margin provides useful evidence
for subsequent architectural reasoning.

The measured execution results are added
to the exploration history and supplied
to the INS state, which extracts
execution-grounded observations and
guidance for the next iteration.



\subsubsection{INS: Insight Extraction and Feedback Generation}
\label{app:illustrative_ins}

Following architectural evaluation, the framework
enters the Insight Extraction (INS) state.
INS analyzes the measured performance, hardware
cost, and soft-error vulnerability outcomes of
the selected configuration and generates
structured feedback for subsequent exploration.

Unlike the deterministic candidate-processing
states, INS uses LLM-based reasoning to interpret
the execution results, identify remaining
weaknesses, and recommend architectural
adjustments.

In this illustrative walkthrough, the selected
candidate C1 corresponds to a configuration
evaluated in an actual REMODEL-V execution.
Therefore, its recorded INS response is reused
to demonstrate the feedback-generation process.

\paragraph{Recorded INS Response}

Listing~\ref{lst:illustrative_ins}
presents the original INS response generated
after evaluating C1.

\begin{lstlisting}[
    style=remodeljson,
    caption={Recorded INS feedback generated
    after evaluating candidate C1.},
    label={lst:illustrative_ins}
]
{
  "iteration": 1,
  "performance_trend": "The executed design demonstrates a similar performance trend compared to previous designs, showing competitive IPC and hit rates, indicating a strong alignment with the performance metrics.",
  "main_observations": [
    "The IPC of the executed design is 1.15, which is satisfactory and indicates moderate performance under the current configuration.",
    "The hit rates for both the L1 data cache and the MSHRs are extremely high (0.9999 and ~0.9973), suggesting efficient cache utilization.",
    "Vulnerability metrics, particularly vulnerability per tick for L2 MSHR (0.00486) and dcache MSHR (0.0958), indicate vulnerability levels within acceptable limits but should be watched for future designs."
  ],
  "remaining_weaknesses": [
    "The dynamic power (2.272W) is close to its upper bound (2.3W), making it a potential risk area. A reduction in dynamic power is necessary for future iterations.",
    "The peak power (8.764W) is also nearing its upper limit (9W) and should be addressed to avoid exceeding cost constraints."
  ],
  "good_patterns": [
    "Utilizing a medium-sized L1D cache (16kB) has resulted in strong performance metrics, indicating that this configuration balances performance with cost effectively.",
    "A moderately sized L2 cache (256kB) with 4-way associativity has contributed to good hit rates, suggesting effective cache management."
  ],
  "risky_patterns": [
    "The current configuration of L1D MSHRs (6) and ROB entries (128) may lead to increased dynamic power consumption, warranting careful tuning in future designs to prevent power excess.",
    "The peak power consumption is high, which increases the risk of exceeding cost constraints in future designs."
  ],
  "guidance_for_next_acg": [
    "Focus on configurations that reduce dynamic power, such as decreasing the L1D MSHRs or ROB entries.",
    "Consider re-evaluating the L2 cache size to maintain performance while addressing peak power limits.",
    "Explore alternative configurations for the L1D size that could balance performance with a reduction in power consumption."
  ]
}
\end{lstlisting}

\paragraph{Interpretation of INS Feedback}

Listing~\ref{lst:illustrative_ins}
demonstrates how INS transforms execution
measurements into structured guidance.

\begin{enumerate}

\item \textbf{Performance observations:}
The \texttt{performance\_trend} and
\texttt{main\_observations} fields summarize
the measured IPC, cache hit rate, MSHR hit
rate, and vulnerability values.

The observations identify favorable cache
behavior while acknowledging that
vulnerability remains an optimization
consideration.

\item \textbf{Remaining weaknesses:}
The \texttt{remaining\_weaknesses} field
identifies cost metrics approaching their
upper limits.

In particular, dynamic power reaches
2.272 W against an upper bound of 2.3 W,
while peak power reaches 8.764 W against
an upper bound of 9.0 W.

Although both values satisfy the cost
constraints, the remaining margins
are relatively small.

\item \textbf{Good architectural patterns:}
The \texttt{good\_patterns} field identifies
architectural choices that the LLM
associates with favorable observed
performance.

In this example, INS highlights the
16 kB L1D cache and 256 kB L2 cache
with four-way associativity.

\item \textbf{Risky architectural patterns:}
The \texttt{risky\_patterns} field identifies
potential architectural risks.

INS suggests that six L1D MSHRs and
128 ROB entries may contribute to
higher dynamic power consumption.

These statements represent architectural
hypotheses inferred by the LLM, rather
than independently verified causal
relationships.

\item \textbf{Guidance for subsequent ACG:}
The \texttt{guidance\_for\_next\_acg}
field converts the observations into
actionable recommendations.

INS recommends investigating smaller
MSHR or ROB configurations, reconsidering
L2 cache capacity, and exploring
alternative L1D cache sizes while
preserving cost feasibility.

\end{enumerate}

The INS response is incorporated into
the execution history and subsequently
used by REF to construct the reasoning
context for iteration 2.

Importantly, the measured configuration
satisfies all seven cost constraints.
Consequently, the framework has now
established a feasible architectural
design, allowing Mode-Guided Reasoning
to transition from Cost-Recovery Mode
to Observable-Improvement Mode.

The next iteration therefore shifts
its reasoning priority toward improving
performance and reducing soft-error
vulnerability while maintaining the
established cost feasibility.


\subsection{Iteration 2: Feedback-Guided Exploration}
\label{app:illustrative_iteration2}

The second iteration demonstrates how REMODEL-V
uses previously observed execution evidence
to guide subsequent architectural exploration.

In iteration 1, the selected configuration C1
satisfied all seven Medium-regime cost constraints.
The resulting measurements and INS feedback
are incorporated into the execution history
and made available to REF.

Because a cost-feasible configuration has now
been observed, the Mode-Guided Reasoning (MGR)
policy transitions from Cost-Recovery Mode
to Observable-Improvement Mode.

Under this mode, the framework prioritizes
improving the observed performance and
soft-error vulnerability objectives while
preserving cost feasibility.

\subsubsection{REF: Updated Reference Context}
\label{app:illustrative_ref2}

At the beginning of iteration 2, REF constructs
an updated reasoning context using the previously
executed configuration, its measured architectural
outcomes, cost-feasibility status, and INS feedback.

Unlike the first iteration, REF now contains
execution-grounded evidence that can be used
to evaluate architectural trade-offs and
guide new candidate generation.

\paragraph{Recorded REF Context}

Listing~\ref{lst:illustrative_ref2}
presents a condensed excerpt of the actual
REF context recorded at the beginning of
iteration 2.

The excerpt preserves the previously executed
configuration, selected measured outcomes,
cost-feasibility status, and INS guidance.
The complete REF record contains additional
ranking information and feedback fields.

\begin{lstlisting}[
    style=remodeljson,
    caption={Condensed REF context at the beginning
    of iteration 2, based on recorded execution data.},
    label={lst:illustrative_ref2}
]
{
  "iteration": 2,
  "previously_executed_designs": [
    {
      "l1d_size": "16kB",
      "l1d_mshrs": 6,
      "rob_entries": 128,
      "lq_entries": 16,
      "l2_size": "256kB",
      "l2_assoc": 4
    }
  ],
  "ranked_executed_designs": [
    {
      "iteration": 1,
      "cost_feasible": true,
      "metrics_used_for_ranking": {
        "ipc": 1.15112703986,
        "hit_rate": 0.999986879309,
        "mshr_hit_rate": 0.997260273973,
        "access_throughput": 326404710.921,
        "system.l2.mshr.vulnerability_per_tick": 0.004863458035739129,
        "system.cpu.dcache.mshr.vulnerability_per_tick": 0.09580629212425008
      },
      "observable_score": 0.5,
      "rank": 1
    }
  ],
  "remaining_weaknesses": [
    "The dynamic power (2.272W) is close to its upper bound (2.3W), making it a potential risk area.",
    "The peak power (8.764W) is also nearing its upper limit (9W)."
  ],
  "guidance_for_next_acg": [
    "Prioritize cost feasibility first.",
    "Among cost-feasible designs, prefer higher observable_score.",
    "Do not repeat previously executed designs.",
    "Focus on configurations that reduce dynamic power, such as decreasing the L1D MSHRs or ROB entries.",
    "Consider re-evaluating the L2 cache size to maintain performance while addressing peak power limits.",
    "Explore alternative configurations for the L1D size that could balance performance with a reduction in power consumption."
  ]
}
\end{lstlisting}

\paragraph{Interpretation of the Updated REF Context}

Listing~\ref{lst:illustrative_ref2}
demonstrates how REF incorporates information
from the preceding execution.

\begin{enumerate}

\item \textbf{Previously executed configuration:}
The \texttt{previously\_executed\_designs}
field records the configuration selected
during iteration 1.

This information allows subsequent
candidate-processing states to identify
and avoid configurations that have
already been evaluated.

\item \textbf{Measured architectural outcomes:}
The \texttt{ranked\_executed\_designs}
field contains the measured performance
and vulnerability outcomes of the
previously executed configuration.

These values provide an observed reference
for subsequent architectural reasoning.

\item \textbf{Cost-feasibility status:}
The recorded configuration is marked
\texttt{cost\_feasible: true} because
all seven measured cost values satisfy
the active Medium-regime constraints.

This establishes that a feasible
configuration has been discovered.

\item \textbf{Remaining weaknesses:}
REF incorporates the concerns identified
by INS, particularly the limited remaining
margins in dynamic and peak power.

\item \textbf{Guidance for ACG:}
The updated guidance encourages ACG
to investigate alternative architectural
configurations that may reduce hardware
cost while preserving or improving
performance and vulnerability outcomes.

\end{enumerate}

\paragraph{Transition to Observable-Improvement Mode}

The active reasoning mode is determined
using the feasibility information
accumulated from executed configurations.

At the beginning of iteration 1, no
feasible configuration was available.
Therefore, Cost-Recovery Mode was active.

After iteration 1, the executed
configuration satisfies all cost bounds.
Consequently, Observable-Improvement Mode
becomes active for iteration 2.

This transition changes the reasoning
priority from establishing feasibility
to improving the observed objectives
while maintaining feasibility.

Importantly, the transition is based
on measured cost-feasibility information,
not on the LLM's prediction of whether
a candidate is feasible.

The updated REF context and active
reasoning instructions are subsequently
provided to the ACG instances for
the generation of new architectural
proposals.


\subsubsection{ACG: Feedback-Guided Candidate Generation}
\label{app:illustrative_acg2}

Following the updated REF context, REMODEL-V
enters the Architectural Candidate Generation
(ACG) state for the second iteration.

Unlike the initial iteration, ACG now receives
execution-grounded information about the
previously evaluated configuration, including
its architectural parameters, measured outcomes,
cost-feasibility status, and INS feedback.

Because a feasible configuration has already
been discovered, Observable-Improvement Mode
is active. Under this mode, ACG is guided
toward improving the observed performance
and vulnerability objectives while preserving
cost feasibility.

To illustrate how previous execution evidence
influences candidate generation,
Listing~\ref{lst:illustrative_acg2}
presents the actual LLM reasoning recorded
for ACG\_0 during iteration 2.

\paragraph{Recorded LLM Response for ACG\_0}

\begin{lstlisting}[
    style=remodeljson,
    caption={Recorded ACG\_0 response demonstrating
    feedback-guided architectural reasoning
    during iteration 2.},
    label={lst:illustrative_acg2}
]
{
  "candidate_id": "ACG_0",
  "params": {
    "l1d_size": "16kB",
    "l1d_mshrs": 4,
    "rob_entries": 64,
    "lq_entries": 16,
    "l2_size": "512kB",
    "l2_assoc": 4
  },
  "trace": {
    "retrospect": "The previous design had a good observable score but was close to violating dynamic and peak power limits, requiring caution in future configurations.",
    "analysis": {
      "steps": [
        {
          "thinking": "Reducing 'l1d_mshrs' from 6 to 4 may decrease dynamic power consumption while maintaining reasonable performance.",
          "step_conclusion": "Lowering 'l1d_mshrs' potentially reduces dynamic power."
        },
        {
          "thinking": "Reducing 'rob_entries' from 128 to 64 could help in reducing overall power consumption without severely impacting performance.",
          "step_conclusion": "Lowering 'rob_entries' may help decrease both dynamic and peak power."
        },
        {
          "thinking": "Increasing 'l2_size' from '256kB' to '512kB' can improve cache performance, potentially enhancing IPC and hit rates without significantly impacting cost.",
          "step_conclusion": "Increasing 'l2_size' might enhance performance metrics."
        }
      ],
      "final_conclusion": "This configuration balances performance with reduced power consumption, enhancing cost feasibility."
    },
    "explanation": "This candidate reduces 'l1d_mshrs' and 'rob_entries' to lower the dynamic power while increasing 'l2_size' to enhance performance metrics, making it a balanced design."
  }
}
\end{lstlisting}

\paragraph{Interpretation of Feedback-Guided Reasoning}

Listing~\ref{lst:illustrative_acg2}
demonstrates how execution-grounded
feedback influences the LLM's
architectural reasoning.

\begin{enumerate}

\item \textbf{Architectural configuration
(Lines 2--10):}
ACG\_0 proposes a configuration containing
16~kB L1D capacity, four L1D MSHRs,
64 ROB entries, 16 LQ entries, a
512~kB L2 cache, and four-way
L2 associativity.

\item \textbf{Retrospective reasoning
(Line 12):}
The \texttt{retrospect} field explicitly
references the previously evaluated
configuration and the observed proximity
of dynamic and peak power to their
upper bounds.

This demonstrates that ACG now reasons
using previous execution evidence rather
than only the initial user requirements.

\item \textbf{Reducing L1D MSHRs
(Lines 15--18):}
The LLM proposes reducing the number
of L1D MSHRs from six to four.

It hypothesizes that this adjustment
may reduce dynamic power consumption
while maintaining acceptable performance.

\item \textbf{Reducing ROB capacity
(Lines 19--22):}
The LLM proposes reducing ROB entries
from 128 to 64.

Its reasoning suggests that a smaller
ROB may reduce dynamic and peak power,
although the actual effects must be
determined through execution.

\item \textbf{Increasing L2 capacity
(Lines 23--26):}
The LLM proposes increasing L2 cache
capacity from 256~kB to 512~kB.

This change is intended to improve
cache-related performance while
balancing the cost reductions expected
from the other parameter adjustments.

\item \textbf{Overall conclusion and
explanation (Lines 28--30):}
The final fields summarize the LLM's
expectation that the combined parameter
changes will balance performance and
power consumption while maintaining
cost feasibility.

\end{enumerate}

The generated configuration demonstrates
how REMODEL-V uses INS feedback to guide
architectural changes across iterations.

Importantly, the reasoning represents
predicted architectural effects rather
than verified causal relationships.
Because multiple parameters are modified
simultaneously, subsequent measurements
cannot independently establish the
contribution of each parameter change.

The proposed configuration is forwarded
to the deterministic candidate-processing
states before final selection and execution.



\paragraph{Candidate Pool for Iteration 2}
\label{app:illustrative_acg2_pool}

To illustrate candidate processing under
Observable-Improvement Mode, the worked example
considers six architectural proposals.

The first three configurations reproduce the
ACG outputs recorded during the second iteration
of the reference execution. Three additional
configurations are constructed to demonstrate
how REMODEL-V handles repeated architectural
designs when execution history is available.

Table~\ref{tab:illustrative_acg2_pool}
summarizes the six proposals.

\begin{table}[htbp]
\centering
\caption{Architectural candidate pool used in
the illustrative second iteration.}
\label{tab:illustrative_acg2_pool}
\scriptsize
\setlength{\tabcolsep}{3pt}
\renewcommand{\arraystretch}{1.1}
\begin{tabular}{lcccccc}
\hline
\textbf{Parameter} &
\textbf{ACG\_0} &
\textbf{ACG\_1} &
\textbf{ACG\_2} &
\textbf{ACG\_3} &
\textbf{ACG\_4} &
\textbf{ACG\_5} \\
\hline
L1D size (kB) & 16 & 32 & 8 & 16 & 16 & 32 \\
L1D MSHRs & 4 & 4 & 4 & 6 & 4 & 6 \\
ROB entries & 64 & 128 & 64 & 128 & 64 & 64 \\
LQ entries & 16 & 16 & 16 & 16 & 16 & 16 \\
L2 size (kB) & 512 & 512 & 512 & 256 & 512 & 256 \\
L2 associativity & 4 & 4 & 4 & 4 & 4 & 8 \\
\hline
\end{tabular}
\end{table}

The first three candidates reflect different
architectural adjustments proposed using the
updated REF context and Observable-Improvement
Mode instructions.

ACG\_0 reduces L1D MSHRs and ROB capacity
while increasing L2 cache size, consistent
with the recorded reasoning shown in
Listing~\ref{lst:illustrative_acg2}.

ACG\_1 and ACG\_2 provide alternative
architectural resource allocations.

The remaining three proposals are illustrative:

\begin{itemize}

\item \textbf{ACG\_3:}
Repeats the configuration executed in
iteration 1. This candidate illustrates
how the framework detects a previously
evaluated architectural design.

\item \textbf{ACG\_4:}
Duplicates ACG\_0 within the current
candidate pool, illustrating
within-iteration duplicate detection.

\item \textbf{ACG\_5:}
Provides an additional architecturally
valid and unique configuration.

\end{itemize}

All six proposals contain parameter values
within the permitted architectural domains.
Consequently, the illustrative NOR and FV
stages retain all six candidates.

During CDR/DED, however, ACG\_3 is identified
as a previously executed configuration,
while ACG\_4 duplicates another proposal
within the current candidate pool.

Both are removed, leaving four unique,
previously unexecuted configurations
for subsequent PRE assessment.

This example demonstrates how execution
history influences candidate processing
in addition to LLM-guided architectural
reasoning.


\subsubsection{NOR, FV, and CDR/DED: Candidate Processing}
\label{app:illustrative_processing2}

Following candidate generation, the six
architectural proposals enter the deterministic
candidate-processing stages: Normalization (NOR),
Feasibility Verification (FV), and
Canonicalization and Deduplication (CDR/DED).

These states perform the same operations
described in iteration 1. However, the
second iteration introduces an additional
consideration: the framework now has a
previously executed architectural configuration
that must not be evaluated again.

\paragraph{Normalization and Architectural Validation}

NOR standardizes the six generated parameter
assignments into the required internal
representation.

All six illustrative candidates contain
complete parameter dictionaries and are
successfully normalized without modification.

FV then checks the proposed parameter values
against the architectural domains defined in
Table~\ref{tab:appendix_design_space}.

Because all six candidates contain permitted
parameter values, no candidate is rejected
during architectural validity checking.

Thus,

\begin{equation}
N_{\mathrm{ACG}}
=
N_{\mathrm{NOR}}
=
N_{\mathrm{FV}}
=
6.
\end{equation}

\paragraph{History-Aware Duplicate Elimination}

After FV, CDR/DED constructs a canonical
representation for each candidate and
compares it against both the current
candidate pool and the previously executed
design history.

Listing~\ref{lst:illustrative_cdr2}
presents the illustrative duplicate-detection
outcomes for iteration 2.

\begin{lstlisting}[
    style=remodeljson,
    caption={Illustrative CDR/DED processing
    using current-pool and execution-history
    duplicate detection in iteration 2.},
    label={lst:illustrative_cdr2}
]
[CDR/DED] Iteration 2

ACG_0: KEPT
  Unique candidate.
  Assigned ID: C0

ACG_1: KEPT
  Unique candidate.
  Assigned ID: C1

ACG_2: KEPT
  Unique candidate.
  Assigned ID: C2

ACG_3: REJECTED
  Reason: Previously executed configuration.
  Matches iteration 1 selected design.

ACG_4: REJECTED
  Reason: Duplicate in current candidate pool.
  Identical to ACG_0.

ACG_5: KEPT
  Unique candidate.
  Assigned ID: C3

Input candidates: 6
Previously executed duplicates: 1
Current-pool duplicates: 1
Unique candidates retained: 4
\end{lstlisting}

\paragraph{Interpretation of Candidate Processing}

Listing~\ref{lst:illustrative_cdr2}
demonstrates two distinct duplicate-detection
mechanisms.

\begin{enumerate}

\item \textbf{Unique configurations:}
ACG\_0, ACG\_1, ACG\_2, and ACG\_5
contain distinct architectural parameter
assignments that have not been executed
previously. These four candidates are
retained for subsequent assessment.

\item \textbf{Previously executed duplicate:}
ACG\_3 reproduces the configuration
selected during iteration 1:

\begin{equation}
x_{\mathrm{ACG}_3}^{(2)}
=
x_{\mathrm{C1}}^{(1)}.
\end{equation}

Because the configuration has already
been evaluated, CDR/DED removes it
before another architectural evaluation
can occur.

\item \textbf{Within-pool duplicate:}
ACG\_4 has the same architectural
parameter assignment as ACG\_0:

\begin{equation}
K(x_{\mathrm{ACG}_4}^{(2)})
=
K(x_{\mathrm{ACG}_0}^{(2)}).
\end{equation}

Therefore, ACG\_4 is removed as a
duplicate within the current
candidate pool.

\item \textbf{Final candidate pool:}
After eliminating both duplicates,
four unique configurations remain.

\end{enumerate}

The number of candidates is therefore
reduced from six to four:

\begin{equation}
N_{\mathrm{FV}} = 6,
\qquad
N_{\mathrm{CDR/DED}} = 4.
\end{equation}

Table~\ref{tab:illustrative_cdr2_mapping}
summarizes the candidate-processing results.

\begin{table}[htbp]
\centering
\caption{Illustrative candidate-processing
outcomes during iteration 2.}
\label{tab:illustrative_cdr2_mapping}
\renewcommand{\arraystretch}{1.1}
\begin{tabular}{lcc}
\hline
\textbf{ACG Candidate} &
\textbf{Internal ID} &
\textbf{Outcome} \\
\hline
ACG\_0 & C0 & Retained \\
ACG\_1 & C1 & Retained \\
ACG\_2 & C2 & Retained \\
ACG\_3 & -- & Previously executed \\
ACG\_4 & -- & Within-pool duplicate \\
ACG\_5 & C3 & Retained \\
\hline
\end{tabular}
\end{table}

Unlike the first iteration, where the
execution history was empty, the
second iteration demonstrates how
CDR/DED uses accumulated architectural
evidence to prevent repeated evaluations.

The four retained configurations are
forwarded to PRE for assessment using
the updated REF context.



\subsubsection{PRE: Execution-Grounded Candidate Assessment}
\label{app:illustrative_pre2}

Following deterministic candidate processing,
the four retained configurations enter the
Pre-Decision Evaluation (PRE) state.

Unlike the first iteration, PRE now has access
to execution-grounded information from the
previously evaluated configuration.

The REF context contains the measured
performance, soft-error vulnerability,
hardware cost, and feasibility status
of the iteration-1 design, together with
the feedback generated by INS.

PRE uses this information to assess whether
each proposed configuration is likely to
provide a useful architectural improvement
while respecting the active cost constraints.

\paragraph{Illustrative PRE Response}

Listing~\ref{lst:illustrative_pre2}
presents the candidate assessments used in
this worked example.

The assessments for C0, C1, and C2 reproduce
the recorded PRE decisions from the original
three-candidate execution. The assessment
for C3 is constructed to illustrate the
expanded four-candidate pool.

\begin{lstlisting}[
    style=remodeljson,
    caption={Illustrative PRE assessment
    of four candidates in iteration 2.},
    label={lst:illustrative_pre2}
]
[
  {
    "candidate_id": "C0",
    "source_candidate_id": "ACG_0",
    "send_to_DEC": true,
    "confidence": 0.8,
    "reason": "Candidate C0 has lower L1D MSHRs and ROB entries, which may help reduce dynamic power and peak power, making it more likely to remain cost-feasible."
  },
  {
    "candidate_id": "C1",
    "source_candidate_id": "ACG_1",
    "send_to_DEC": false,
    "confidence": 0.5,
    "reason": "Candidate C1 retains a higher number of ROB entries, which may increase dynamic power significantly, making it less likely to comply with cost constraints."
  },
  {
    "candidate_id": "C2",
    "source_candidate_id": "ACG_2",
    "send_to_DEC": true,
    "confidence": 0.7,
    "reason": "Candidate C2 has the lowest L1D size, which could lead to lower power consumption and maintain acceptable performance, aligning with the guidance to focus on feasibility."
  },
  {
    "candidate_id": "C3",
    "source_candidate_id": "ACG_5",
    "send_to_DEC": false,
    "confidence": 0.6,
    "reason": "The proposed cache and execution resources may increase power consumption and reduce the available cost margin."
  }
]
\end{lstlisting}

\paragraph{Interpretation of PRE Assessments}

Listing~\ref{lst:illustrative_pre2}
demonstrates how PRE evaluates candidates
using previous execution evidence.

\begin{enumerate}

\item \textbf{Candidate C0: Forwarded to DEC.}
PRE predicts that reducing L1D MSHRs
and ROB capacity may reduce dynamic
and peak power consumption.

The candidate is forwarded with
a recorded confidence value of 0.8.

\item \textbf{Candidate C1: Rejected by PRE.}
PRE identifies the relatively large
ROB capacity as a potential source
of increased dynamic power.

Because the candidate is considered
less likely to preserve cost feasibility,
PRE excludes it from the final
decision-making pool.

\item \textbf{Candidate C2: Forwarded to DEC.}
PRE considers the smaller L1D cache
a potentially useful adjustment
for reducing power consumption.

The candidate is forwarded with
a recorded confidence value of 0.7.

\item \textbf{Candidate C3: Illustrative rejection.}
The additional constructed candidate
is excluded because its architectural
resource allocation is assessed as
potentially increasing power consumption.

This assessment is included for
explanatory purposes and was not
generated by PRE during the original
recorded experiment.

\end{enumerate}

The illustrative PRE process reduces
the candidate pool from four
configurations to two:

\begin{equation}
N_{\mathrm{CDR/DED}} = 4,
\qquad
N_{\mathrm{PRE}} = 2.
\end{equation}

The resulting candidate pool contains
C0 and C2, which are forwarded to DEC
for final architectural selection.

Importantly, PRE's assessments are
based on LLM reasoning over previously
observed evidence. The forwarding
decisions do not establish the actual
performance or cost feasibility of
unexecuted candidates.

The final decision remains the
responsibility of DEC, while EXC
provides the measured architectural
outcomes needed to evaluate the
selected configuration.



\subsubsection{DEC: Feedback-Guided Candidate Selection}
\label{app:illustrative_dec2}

Following PRE assessment, the framework enters
the Decision (DEC) state. In the illustrative
second iteration, PRE forwards two candidates,
C0 and C2, for final architectural selection.

Unlike the first iteration, DEC now receives
execution-grounded information from the
previously evaluated configuration, including
its measured performance, vulnerability,
hardware cost, and feasibility status.

The active reasoning policy is
Observable-Improvement Mode. DEC must therefore
consider potential performance and vulnerability
improvements while preserving the established
cost feasibility.

\paragraph{Candidate Pool Supplied to DEC}

Table~\ref{tab:illustrative_dec2_candidates}
summarizes the two candidates forwarded by PRE.

\begin{table}[htbp]
\centering
\caption{Candidates forwarded to DEC during
the illustrative second iteration.}
\label{tab:illustrative_dec2_candidates}
\renewcommand{\arraystretch}{1.1}
\begin{tabular}{lcc}
\hline
\textbf{Parameter} &
\textbf{C0} &
\textbf{C2} \\
\hline
L1D size & 16 kB & 8 kB \\
L1D MSHRs & 4 & 4 \\
ROB entries & 64 & 64 \\
LQ entries & 16 & 16 \\
L2 size & 512 kB & 512 kB \\
L2 associativity & 4 & 4 \\
\hline
\end{tabular}
\end{table}

Both candidates reduce L1D MSHRs and ROB
capacity relative to the configuration
executed during iteration 1.

The principal architectural difference
between C0 and C2 is the L1D cache size:
C0 retains 16 kB, whereas C2 proposes
8 kB.

These configurations are the same two
candidates supplied to DEC in the
recorded second iteration.

\paragraph{Recorded DEC Response}

Listing~\ref{lst:illustrative_dec2}
presents the actual DEC response
generated for this candidate pool.

\begin{lstlisting}[
    style=remodeljson,
    caption={Recorded DEC response selecting
    candidate C0 during iteration 2.},
    label={lst:illustrative_dec2}
]
{
  "chosen_id": "C0",
  "explanation": "Candidate C0 is expected to provide a balanced approach to cost-feasibility and performance due to its conservative settings, which target the reduction of dynamic and peak power."
}
\end{lstlisting}

\paragraph{Interpretation of the DEC Response}

The recorded DEC response contains two
principal fields.

\begin{enumerate}

\item \textbf{Selected candidate (Line 2):}
The \texttt{chosen\_id} field identifies
C0 as the architectural configuration
selected for execution.

DEC chooses an existing candidate from
the PRE-approved pool rather than
generating a new architectural design.

\item \textbf{Selection explanation (Line 3):}
The \texttt{explanation} field states
that C0 is expected to balance cost
feasibility and performance through
conservative architectural settings.

In particular, the LLM emphasizes
the expected reduction in dynamic
and peak power consumption.

\end{enumerate}

The selected configuration is

\begin{equation}
x_{\mathrm{C0}}^{(2)}
=
(
16\text{ kB},
4,
64,
16,
512\text{ kB},
4
),
\label{eq:illustrative_dec2_design}
\end{equation}

where the parameter ordering is
L1D size, L1D MSHRs, ROB entries,
LQ entries, L2 size, and L2 associativity.

Relative to the previously executed
configuration, the selected design
reduces L1D MSHRs from six to four
and ROB entries from 128 to 64,
while increasing L2 cache capacity
from 256 kB to 512 kB.

These modifications are consistent
with the architectural guidance
generated by INS after iteration 1.

However, the DEC explanation remains
a prediction based on the available
execution evidence. The actual
performance, vulnerability, and
cost effects of the selected design
must be determined through EXC.

The selected configuration is therefore
forwarded to EXC for architectural
evaluation.



\subsubsection{EXC: Architectural Execution and Measurement}
\label{app:illustrative_exc2}

Following DEC selection, the framework enters
the Execution (EXC) state to obtain the
architectural measurements of candidate C0.

The selected configuration reduces L1D MSHRs
from six to four and ROB entries from 128
to 64, while increasing L2 cache capacity
from 256 kB to 512 kB.

These modifications were proposed to reduce
hardware power consumption while maintaining
favorable architectural performance.

As in iteration 1, the database-backed
execution mode retrieves previously measured
architectural outcomes for the selected
configuration. No new gem5 simulation is
launched during this illustrative execution.

\paragraph{Recorded Execution Output}

Listing~\ref{lst:illustrative_exc2}
presents the measured outcomes of the
selected configuration.

The listing is a condensed representation
of the recorded execution result, preserving
the original architectural measurements.

\begin{lstlisting}[
    style=remodeljson,
    caption={Recorded architectural measurements
    for candidate C0 during iteration 2.},
    label={lst:illustrative_exc2}
]
{
  "status": "evaluated_from_db",
  "design": {
    "l1d_size": "16kB",
    "l1d_mshrs": 4,
    "rob_entries": 64,
    "lq_entries": 16,
    "l2_size": "512kB",
    "l2_assoc": 4
  },
  "metrics": {
    "ipc": 0.898868536284,
    "hit_rate": 0.999986879654,
    "mshr_hit_rate": 0.997138769671,
    "access_throughput": 254908026.756,
    "system.l2.mshr.vulnerability_per_tick": 0.003654447015448587,
    "system.cpu.dcache.mshr.vulnerability_per_tick": 0.07381016964449114,
    "total_power": 2.25449602,
    "area": 4.27107,
    "dynamic_power": 2.18413,
    "peak_power": 8.61839602,
    "leakage_power": 0.07036602,
    "total_energy": 0.00134818861996,
    "edp": 8.06216794736e-07
  }
}
\end{lstlisting}

\paragraph{Measured Cost Feasibility}

The retrieved cost measurements are compared
against the same Medium-regime constraints
used during iteration 1.

Table~\ref{tab:illustrative_exc2_cost}
summarizes the resulting feasibility assessment.

\begin{table}[htbp]
\centering
\caption{Measured cost feasibility of
candidate C0 during iteration 2.}
\label{tab:illustrative_exc2_cost}
\small
\renewcommand{\arraystretch}{1.1}
\begin{tabular}{lccc}
\hline
\textbf{Cost Metric} &
\textbf{Measured} &
\textbf{Upper Bound} &
\textbf{Status} \\
\hline
Total power & 2.2545 & 2.5 & Feasible \\
Area & 4.2711 & 5.0 & Feasible \\
Dynamic power & 2.1841 & 2.3 & Feasible \\
Peak power & 8.6184 & 9.0 & Feasible \\
Leakage power & 0.070366 & 0.15 & Feasible \\
Total energy & 0.001348 & 0.15 & Feasible \\
EDP & $8.062\times10^{-7}$ & 8.0 & Feasible \\
\hline
\end{tabular}
\end{table}

All seven measured cost values satisfy
their corresponding upper bounds.
Therefore, candidate C0 is classified
as cost-feasible.

\paragraph{Comparison with the Previous Execution}

To determine how the architectural changes
affected the observed objectives,
Table~\ref{tab:illustrative_execution_comparison}
compares the measurements from iterations
1 and 2.

\begin{table}[htbp]
\centering
\caption{Measured architectural outcomes
across the two illustrative iterations.}
\label{tab:illustrative_execution_comparison}
\small
\renewcommand{\arraystretch}{1.1}
\begin{tabular}{lcc}
\hline
\textbf{Metric} &
\textbf{Iteration 1} &
\textbf{Iteration 2} \\
\hline
IPC & 1.1511 & 0.8989 \\
Cache hit rate & 0.999987 & 0.999987 \\
MSHR hit rate & 0.997260 & 0.997139 \\
Access throughput & $3.264\times10^8$ & $2.549\times10^8$ \\
L2-MSHR vulnerability & 0.004863 & 0.003654 \\
L1D-MSHR vulnerability & 0.095806 & 0.073810 \\
Total power & 2.3424 & 2.2545 \\
Dynamic power & 2.2721 & 2.1841 \\
Peak power & 8.7644 & 8.6184 \\
Area & 4.3912 & 4.2711 \\
\hline
\end{tabular}
\end{table}

The measured results show that the second
configuration reduces total, dynamic,
and peak power while preserving cost
feasibility.

Both vulnerability measurements also
decrease, indicating lower measured
soft-error vulnerability for the
selected structures.

However, these improvements are accompanied
by reduced IPC and access throughput.
The second configuration therefore
illustrates a trade-off between
performance, vulnerability, and
hardware cost.

\paragraph{Evaluation of DEC's Expectations}

During candidate selection, DEC expected
the reduced architectural resource
allocation to lower dynamic and peak
power consumption.

The measured outcomes are consistent
with this expectation: dynamic power
decreases from 2.2721 W to 2.1841 W,
while peak power decreases from
8.7644 W to 8.6184 W.

Nevertheless, the experiment changes
multiple architectural parameters
simultaneously. Consequently, the
observed reductions cannot be attributed
independently to the MSHR or ROB changes.

Furthermore, the reduction in IPC
and access throughput demonstrates
that preserving cost feasibility
does not guarantee improvement in
every performance objective.

These execution results are incorporated
into the exploration history and supplied
to INS for feedback generation.



\subsubsection{INS: Execution-Grounded Insight Generation}
\label{app:illustrative_ins2}

Following architectural evaluation, the framework
enters the Insight Extraction (INS) state.
INS analyzes the measured outcomes of the
selected configuration and generates structured
feedback for subsequent exploration.

During iteration 2, the selected configuration
reduces total, dynamic, and peak power, as well
as both soft-error vulnerability measurements.
However, IPC and access throughput decrease
relative to the previously executed design.

These measurements provide the evidence
available to INS for identifying architectural
strengths, remaining weaknesses, and potential
directions for future exploration.

\paragraph{Recorded INS Response}

Listing~\ref{lst:illustrative_ins2} presents
the actual INS response generated after
evaluating the second configuration.

The response is reproduced from the recorded
REMODEL-V execution and is not constructed
for the illustrative six-candidate pool.

\begin{lstlisting}[
    style=remodeljson,
    caption={Recorded INS feedback following
    the second architectural evaluation.},
    label={lst:illustrative_ins2}
]
{
  "iteration": 2,
  "performance_trend": "The executed design demonstrates a robust performance profile, with IPC (0.899) and hit rates (close to 1) indicating strong efficiency in instruction execution and cache utilization. Compared to potential configurations, it performs well but leaves room for improvement in balancing cost and vulnerability metrics.",
  "main_observations": [
    "The executed design met performance goals with high hit rates and acceptable IPC, suggesting the current configuration of L1D size and MSHRs is effective for this application.",
    "Cost metrics such as total power (2.254) and dynamic power (2.184) are near their upper bounds, indicating a risk of violating cost constraints if parameters are aggressively modified in future designs.",
    "Vulnerability metrics for both L2 MSHR (0.0036) and CPU DCache MSHR (0.0738) indicate moderate risk, although they remain within acceptable limits."
  ],
  "remaining_weaknesses": [
    "Total power is close to the upper limit (2.5), indicating future designs should aim to reduce power consumption.",
    "Dynamic power is also near its upper limit (2.3), necessitating careful parameter adjustments to avoid exceeding cost constraints in subsequent iterations."
  ],
  "good_patterns": [
    "The combination of a moderate L1D size (16kB) and adequate MSHRs (4) effectively supports current performance levels while facilitating cost constraints, suggesting these parameters are balanced well.",
    "Using 16 LQ entries appears to manage latency effectively while staying within bounds for area and power."
  ],
  "risky_patterns": [
    "The choice of peak power (8.618) approaches the maximum allowable level (9), posing a risk if further increases are made to capacity without considering power reduction strategies.",
    "Having a relatively low number of ROB entries (64) might limit the performance under higher loads, but this choice helped maintain cost-effectiveness in the current design."
  ],
  "guidance_for_next_acg": [
    "Focus on reducing dynamic and total power by exploring smaller L1D or L2 configurations, or decreasing MSHRs where feasible to enhance power efficiency.",
    "Consider evaluating designs with higher L1D size or more L1D MSHRs for potential performance boosts while carefully monitoring cost impacts. Prioritize configurations that maintain or enhance performance benchmarks without violating power constraints."
  ]
}
\end{lstlisting}

\paragraph{Interpretation of INS Feedback}

Listing~\ref{lst:illustrative_ins2}
demonstrates how INS converts measured
architectural outcomes into structured
reasoning guidance.

\begin{enumerate}

\item \textbf{Performance assessment:}
The \texttt{performance\_trend} and
\texttt{main\_observations} fields describe
the executed configuration's IPC, cache
behavior, hardware cost, and vulnerability.

Although INS describes the configuration
as having favorable performance, the
measured comparison shows that IPC and
access throughput decreased relative to
iteration 1. The LLM's qualitative
assessment should therefore be interpreted
alongside the actual measurements.

\item \textbf{Remaining cost weaknesses:}
INS identifies total and dynamic power
as continuing concerns.

Although both values decreased from
iteration 1 and remain within the
Medium-regime bounds, dynamic power
still has limited remaining margin
before reaching its upper limit.

\item \textbf{Good architectural patterns:}
The \texttt{good\_patterns} field highlights
the 16 kB L1D cache, four L1D MSHRs,
and 16 LQ entries as potentially useful
architectural choices.

These observations represent the LLM's
interpretation of the executed design
rather than independently established
parameter-level causal relationships.

\item \textbf{Risky architectural patterns:}
The \texttt{risky\_patterns} field identifies
peak power and reduced ROB capacity
as potential concerns.

The LLM suggests that a smaller ROB
may limit performance under heavier
execution demands.

\item \textbf{Guidance for future exploration:}
The \texttt{guidance\_for\_next\_acg}
field recommends further exploration
of cache capacities and MSHR allocations
to balance performance improvements
against hardware cost constraints.

This guidance can be incorporated into
REF during the next iteration, allowing
subsequent ACG instances to generate
proposals using the updated execution
history.

\end{enumerate}

The second iteration therefore completes
another execution-grounded feedback cycle.
The selected configuration remains
cost-feasible, but its measured outcomes
reveal trade-offs among performance,
hardware cost, and vulnerability.

Rather than assuming that the selected
design improves every objective, REMODEL-V
retains the measured outcomes and INS
feedback for subsequent exploration.

This illustrates the iterative nature
of the state machine: architectural
proposals are generated, processed,
selected, evaluated, and interpreted,
with the resulting evidence informing
the next reasoning cycle.


\subsection{Summary of the Illustrative State-Machine Execution}
\label{app:illustrative_summary}

The preceding walkthrough demonstrates the
operation of REMODEL-V across two consecutive
exploration iterations. The example combines
recorded LLM reasoning and architectural
measurements with constructed candidate-processing
scenarios to illustrate the principal
mechanisms of the framework.

\paragraph{Candidate Processing across States}

Table~\ref{tab:illustrative_state_summary}
summarizes the number of candidates retained
at each processing stage.

\begin{table}[htbp]
\centering
\caption{Candidate-processing summary for
the two illustrative REMODEL-V iterations.}
\label{tab:illustrative_state_summary}
\renewcommand{\arraystretch}{1.1}
\begin{tabular}{lcc}
\hline
\textbf{State} &
\textbf{Iteration 1} &
\textbf{Iteration 2} \\
\hline
Active reasoning mode & Cost-Recovery & Observable-Improvement \\
ACG proposals & 6 & 6 \\
After NOR & 6 & 6 \\
After FV & 5 & 6 \\
After CDR/DED & 4 & 4 \\
Forwarded by PRE & 4 & 2 \\
Selected by DEC & 1 & 1 \\
Evaluated by EXC & 1 & 1 \\
INS feedback records & 1 & 1 \\
\hline
\end{tabular}
\end{table}

In the first iteration, FV removes one
illustrative candidate containing an
out-of-domain architectural parameter.
CDR/DED then eliminates another candidate
that duplicates an existing proposal
within the same candidate pool.

Because no previously executed design is
available, PRE forwards all four remaining
candidates to DEC. DEC selects C1, whose
previously measured architectural outcomes
satisfy all seven Medium-regime cost
constraints.

In the second iteration, all six
illustrative proposals satisfy the
architectural parameter domains.
However, CDR/DED removes one candidate
that repeats the previously executed
configuration and another that duplicates
a candidate within the current pool.

PRE then uses the available execution
evidence to assess the four remaining
candidates and forwards two to DEC.
DEC selects C0 for architectural
evaluation.

\paragraph{Mode-Guided Reasoning Transition}

The two iterations illustrate the
transition between the reasoning modes
of REMODEL-V.

During iteration 1, the execution
history is empty and no feasible
configuration has been observed.
Cost-Recovery Mode therefore
prioritizes establishing feasibility.

After the first selected configuration
is evaluated, its measured cost values
satisfy all seven active constraints.
This establishes a feasible
architectural reference.

Consequently, Observable-Improvement
Mode becomes active in iteration 2.
The framework uses the previously
measured outcomes and INS feedback
to guide new architectural proposals
while preserving cost feasibility.

The recorded ACG reasoning demonstrates
this feedback mechanism through proposed
reductions in L1D MSHRs and ROB capacity,
together with an increase in L2 cache size.

\paragraph{Observed Architectural Trade-Offs}

The two selected configurations are
evaluated using previously measured
architectural records from the same
Matrix Multiplication design space.

The second configuration reduces
total, dynamic, and peak power,
as well as both selected soft-error
vulnerability metrics.

However, IPC and access throughput
also decrease, demonstrating that
improvements in hardware cost and
vulnerability may be accompanied
by performance degradation.

Both configurations remain
cost-feasible under the Medium regime.

These outcomes illustrate why
REMODEL-V retains individual
performance, vulnerability, and
cost measurements throughout
exploration rather than relying
exclusively on LLM-generated
architectural expectations.

\paragraph{Key Observations}

The worked example highlights several
important properties of REMODEL-V.

First, ACG generates architectural
proposals and structured reasoning,
while deterministic states control
candidate representation, architectural
validity, and redundancy.

Second, FV and CDR/DED serve distinct
purposes. FV enforces architectural
parameter legality, whereas CDR/DED
prevents repeated configurations from
proceeding to selection and execution.

Third, PRE uses the available execution
history to assess candidates before
DEC selects one configuration for
architectural evaluation.

Fourth, EXC supplies measured
architectural evidence, allowing
the framework to determine actual
cost feasibility and observe
performance--vulnerability trade-offs.

Finally, INS converts execution
observations into structured guidance
that informs subsequent exploration,
while Mode-Guided Reasoning adjusts
the optimization priority according
to the observed feasibility history.

The constructed candidate-rejection
and deduplication examples demonstrate
the intended processing mechanisms;
they are not additional experimental
observations. The recorded LLM
responses and architectural measurements
are included to ground the walkthrough
in actual REMODEL-V behavior.

Overall, the case study illustrates
how the REMODEL-V state machine
combines LLM-based architectural
reasoning with deterministic candidate
processing and execution-grounded
feedback to support iterative
microarchitectural design space
exploration.
